\begin{claim}{F.19}{Proved}{The two-track kernel and the exact remaining source subgroup}
Use the literal base configuration and ordered branch labels of
F.17--18.  Partition its ten labels as
\[
 A=\{0,\infty,1,v_\zeta\},\qquad
 U=\{v_u:u\in\mu _5\setminus\{\zeta\}\},\qquad
 Z=\{a_+,a_-\}.
                                                               \tag{F.19.1}
\]
Let \(P_4,P_8,P_6,P_{10}^{\dagger}\) be the full pure
completions on \(A,A\cup U,A\cup Z,A\cup U\cup Z\), respectively.
The actual forgetting maps form the square
\[
 \begin{array}{ccc}
 P_{10}^{\dagger}&\xrightarrow{\ g\ }&P_6\\
 {\scriptstyle f}\downarrow&&\downarrow{\scriptstyle h}\\
 P_8&\xrightarrow{\ f_0\ }&P_4 .
 \end{array}
 \qquad
 \begin{gathered}
 \mathcal Q=P_8\times_{P_4}P_6,\\
 \pi=(f,g):P_{10}^{\dagger}\twoheadrightarrow\mathcal Q,\\
 D^+=\ker\pi=\ker f\cap\ker g .
 \end{gathered}                                                 \tag{F.19.2}
\]
In particular the displayed map to this specified profinite
fibre product is onto.

For each of the eight pairs \((v_u,a_\pm)\), choose either one
of its two collision points in F.17 and retain its actual
source whisker.  Write \(\nu_{u,\pm}\) for the image of that
source meridian under \(\Phi^\dagger\).
Then
\[
 D^+=\overline{\langle\!\langle
               \nu_{u,\pm}:u\ne\zeta
                         \rangle\!\rangle}_{\,P_{10}^{\dagger}},
 \qquad \nu_{u,\pm}\in K^+ .
                                                               \tag{F.19.3}
\]
The normal closure is taken in the indicated full ambient
group.  This does not replace \(K^+\) by that normal closure.

All four label sets in (F.19.1) are invariant under the
specified involution, and its two new labels are fixed
individually.  Denote the induced actions by \(R\) on each
group where no ambiguity arises.  With
\(\mathcal J=\ker f\) and
\(\mathcal J_4=\ker(h:P_6\to P_4)\), one has
\[
 \begin{gathered}
 1\longrightarrow D^+\longrightarrow\mathcal J
       \xrightarrow{\ g\ }\mathcal J_4\longrightarrow1,\\
 C_{\mathcal J_4}(R)=1,\qquad
 H^1(C_2,D^+)=\{1\}.
 \end{gathered}                                                 \tag{F.19.4}
\]
Here the last expression denotes the continuous nonabelian
cohomology set for this actual action.
More explicitly, the map
\[
 D^+\longrightarrow
     \{\delta\in D^+:\delta R(\delta)=1\},\qquad
                b\longmapsto b^{-1}R(b)
                                                               \tag{F.19.5}
\]
is a homeomorphism.
The centralizer map
\[
          \pi:C^\dagger\xrightarrow{\ \sim\ }\mathcal Q^R
                                                               \tag{F.19.6}
\]
is an isomorphism of profinite groups.

For the original residual group \(\mathcal E\), its unique
branch lift \(\ell\) and the source maps of F.17--18,
\[
 g\ell(u)=\widehat\chi_*(q_0u),\qquad
 b(w)=\ell(u)^{-1}\Phi^\dagger(w)\in D^+
               \quad(qw=u,\ u\in\mathcal E).
                                                               \tag{F.19.7}
\]
Thus \(K^+\subset D^+\), and the exact obstruction map of F.18
has the more precise target
\[
 \mathfrak o_D:\mathcal E\longrightarrow D^+/K^+,\qquad
                 u\longmapsto b(w)K^+ .
                                                               \tag{F.19.8}
\]
It is continuous and independent of the source lift \(w\).
Its distinguished fibre is still exactly the fibre
(F.18.10).

Here are explicit subgroup maps accompanying that comparison.
Set \(H^\dagger=\Phi^\dagger(L^\dagger)\) and
\[
 \begin{aligned}
 \gamma:L&\longrightarrow\mathcal Q,&
       \gamma(u)&=(\widehat\phi_*(u),\widehat\chi_*(q_0u)),\\
 \mathcal T&=\pi^{-1}(\gamma(L)).
 \end{aligned}                                                 \tag{F.19.9}
\]
The map \(\gamma\) is an embedding.  There is a commutative
diagram of exact full profinite sequences
\[
 \begin{array}{ccccccccc}
 1&\longrightarrow&K^+&\longrightarrow&H^\dagger
     &\xrightarrow{\ q_H\ }&L&\longrightarrow&1\\
 &&\downarrow&&\downarrow&&\Vert\\
 1&\longrightarrow&D^+&\longrightarrow&\mathcal T
     &\xrightarrow{\ \pi_L\ }&L&\longrightarrow&1 .
 \end{array}                                                   \tag{F.19.10}
\]
Both rows split by the same specified map
\(j=\Phi^\dagger i\), where \(i:L\to L^\dagger\) retains
the eleven old free meridians in the source filling.
The corresponding inverse coordinates are
\[
 \begin{aligned}
 D^+\rtimes_j L&\longrightarrow\mathcal T,
       &(d,u)&\longmapsto d\,j(u),\\
 \mathcal T&\longrightarrow D^+\rtimes_j L,
       &t&\longmapsto
          \bigl(tj(\pi_Lt)^{-1},\pi_Lt\bigr).
 \end{aligned}                                                 \tag{F.19.11}
\]
Restricting these maps gives \(K^+\rtimes_jL\cong H^\dagger\).
Moreover
\[
                  \mathcal T\cap C^\dagger=\ell(\mathcal E).
                                                               \tag{F.19.12}
\]
These are assertions about the actual images and inverse maps;
they do not assert \(H^\dagger=\mathcal T\) or \(K^+=D^+\).
The source-membership equation in (F.19.8), and triviality
of \(\mathcal E\), remain unresolved.
\end{claim}

\begin{proof}
\emph{The literal six-point centralizer.}
The retained involution on \(A\cup Z\) exchanges
\((0,\infty)\) and \((1,v_\zeta)\), and fixes \(a_+,a_-\).
Let \(C_6^\dagger\) denote its full pure centralizer,
and let \(C_4\) denote the full pure centralizer after
forgetting \(Z\).
There is a branch-permutation character on \(C_4\);
write \(C_4^+\) for its kernel.
The discrete decorated quotient has two ordinary orbit
labels and two branch labels.
The deck involution is nonpure before and after adding
the branch labels, so descent is injective on both
pure centralizers.
The actual inverse lift argument of F.18.2 therefore gives
\[
              h:C_6^\dagger\xrightarrow{\ \sim\ }C_4^+ .
                                                               \tag{F.19.13}
\]
For this application the prime-order completion and full
topology comparison F.7.4--5 is applied to these very
four- and six-point actions.
The branch permutation is onto: exchange the two
quotient branch points, then correct any of the two
ordinary pair-name changes by pushing that ordinary
quotient point once around a branch point.
These corrections preserve the branch permutation.
Thus this is the same exact index-two comparison as
in F.18, with two old pairs retained.

In fact
\[
             \widehat\chi_*(\widehat{\pi_1(Y_0,s_0)})
                                =C_6^\dagger .                \tag{F.19.14}
\]
To check equality with this particular map, descend by
the actual quotient \(Q_a(V)=((V-a)/(V+a))^2\).
F.17.12 identifies its four ordered quotient values
as \(0,\infty,1,t\), \(t=s^2/\zeta^2\).
There are exactly two choices of pure ordinary sheet
names modulo simultaneous deck exchange.
The sheet change \(s\mapsto-s\), followed by
\(V\mapsto V/a^2\), exchanges precisely the names in
\(\{1,a^2\}\); the deck exchanges both ordinary pairs.
Hence the degree-two map \(Y_0\to\mathcal M_{0,4}\)
is exactly the covering belonging to the discrete
pure-lift subgroup of this quotient group.
Its fundamental group is the entire image of the
discrete six-point pure centralizer under descent.
The family of involutions on the six labelled points
shows that \(\chi_*\) lands in that centralizer.
Injective descent now proves the discrete equality.
The full topology of \(\chi_*\), proved using the
ordinary \(a^2\)-cover in F.17, and the prime-order
centralizer comparison complete that equality to
(F.19.14).
Both the quotient parameter \(t\) and ordinary
forgetting parameter \(a^2\) are retained in this
argument; no identification of those two maps is used.

\emph{The actual fibre-product map and its completion.}
Normalize the three common labels \(0,\infty,1\).
The space
\[
 X=\mathcal M(A\cup U)
          \times_{\mathcal M(A)}\mathcal M(A\cup Z)
                                                               \tag{F.19.15}
\]
has coordinates \(v_\zeta\), four ordered \(U\)-coordinates
and two ordered \(Z\)-coordinates.
The \(U\)-coordinates are distinct and avoid \(A\);
the \(Z\)-coordinates are distinct and avoid \(A\).
At this stage a \(U\)-coordinate may equal a
\(Z\)-coordinate.
Both projections to the four-point base are the
actual configuration forgetting fibrations.
The base \(\mathcal M_{0,4}\) has \(\pi_2=0\).
The homotopy sequences therefore identify \(\pi_1(X)\)
with the discrete fibre product of the two pure
mapping class groups over the four-point group:
the kernel on both sides is the product of the
two configuration-fibre groups, and the maps on
those kernels and on the base group are identities.

Here is also the full-topology comparison for this
particular fibre product.
The map \(\mathcal M(A\cup Z)\to\mathcal M(A)\)
has a continuous section.  For the variable coordinate
\(u=v_\zeta\), place the two new coordinates at
\(2+|u|\) and \(3+|u|\).
They are distinct and avoid \(0,1,u,\infty\).
A path in the connected fibre identifies the
resulting basepoint with the specified \((a_+,a_-)\).
This gives a based group section.
The discrete kernel is
\[
 J_{4,\mathrm d}
   =\pi_1(\operatorname{Conf}_2(S_{0,4}),(a_+,a_-)),
 \qquad
 1\longrightarrow F_4\longrightarrow J_{4,\mathrm d}
                     \longrightarrow F_3\longrightarrow1 .
                                                               \tag{F.19.16}
\]
The latter sequence splits by the same far-point
construction in a punctured-plane coordinate.
In particular \(J_{4,\mathrm d}\) is finitely generated.

Consequently the discrete group in (F.19.15) is
\(J_{4,\mathrm d}\rtimes P\Gamma(A\cup U)\), with
the action pulled back from the displayed actual
section over \(A\).
For every finite-index subgroup of the kernel,
choose a characteristic finite-index subgroup inside
it.  The action on this finite quotient has finite
image; quotienting the acting group by its action
kernel gives a finite semidirect quotient of the
whole group.  This proves the full induced topology
on the kernel.  It also extends the action to the
completions and proves exactness of the completed
split sequence: finite homomorphisms on both factors
are separated by these compatible finite quotients.
The completed Birman comparison identifies its kernel
with the actual \(\mathcal J_4\) in \(P_6\).
The same section thus gives
\[
 \widehat{\pi_1(X)}
   \cong\mathcal J_4\rtimes P_8
   \cong P_8\times_{P_4}P_6=\mathcal Q
                                                               \tag{F.19.17}
\]
with the actual projection maps.
This proves completion for the specified fibre product.

For clarity, surjectivity in (F.19.2) can also be
seen directly on those actual full maps.
Lift a given \(P_8\) component through \(f\).
Its \(P_6\) component differs from the desired one
by an element of \(\mathcal J_4\).
The filling map
\[
 \pi_1\operatorname{Conf}_2(S_{0,8})
                  \longrightarrow
       \pi_1\operatorname{Conf}_2(S_{0,4})
\]
is onto.  Apply the two Fadell--Neuwirth sequences:
both the one-point quotient \(F_7\to F_3\) and
the fibre map \(F_8\to F_4\) are the onto puncture
fillings, so a lift on the quotient can be corrected
in the fibre.
Its completion is onto as well, and these completions
are exactly \(\mathcal J\to\mathcal J_4\) by the
completed Birman comparisons.
Correcting the chosen lift inside \(\mathcal J\)
gives the required pair.
This also proves the exact sequence in (F.19.4).

\emph{The eight actual cross-diagonal generators.}
The ten-point moduli space is the open subset of
\(X\) obtained by removing exactly
\[
       \mathcal D_{u,\pm}=\{Z_u=Z_\pm\}\qquad(u\ne\zeta).
                                                               \tag{F.19.18}
\]
These are eight smooth complex hypersurfaces.
Each has connected smooth part after removing its
intersections with the others: identify the coincident
pair, then move the remaining finitely many labelled
coordinates along paths avoiding their finite forbidden
sets.
Different hypersurfaces meet only in complex codimension
at least two; their possible intersections represent
independent pair collisions.

Every loop in \(X\) can be moved off these hypersurfaces.
A disk witnessing a nullhomotopy in \(X\) can be made
transverse to their smooth parts and disjoint from
their intersections.  Its finitely many intersection
points, joined by based arcs in the disk, express
its original boundary as a product of conjugates
of the corresponding signed normal meridians.
Conversely every such meridian dies in \(X\).
This proves that the discrete kernel is the normal
closure of one based meridian for each hypersurface.

Choose its based path to be the actual source
whisker from F.17 approaching either of the two
simple collision points for that pair.
The other labels remain distinct there.
The coefficient-one local calculation (F.17.7--9)
identifies its normal meridian exactly with
\(\nu_{u,\pm}\), with that path and a positive
one-turn exponent.
The finite-target universal property of a quotient
now identifies the kernel on full completions
with the closed normal closure of these eight
elements.  Using (F.19.17) gives exactly (F.19.3).
Every chosen source meridian is killed by \(q\),
so every \(\nu_{u,\pm}\) belongs to the specified
subgroup \(K^+\).

There is a useful exact check on what these ordinary
winding coordinates detect.
The eight global functions \(Z_u-Z_\pm\) give
a continuous homomorphism
\[
 \operatorname{wind}:P_{10}^{\dagger}\longrightarrow
                 \widehat\Z^{\,8},\qquad
 \operatorname{wind}(\nu_{u,\pm})=e_{u,\pm}.
                                                               \tag{F.19.19}
\]
At either root for a given pair, its own coordinate
has winding one and the other seven have winding zero.
The selected eight elements already show that the
images of both \(K^+\) and \(D^+\) under this
particular homomorphism are all of
\(\widehat\Z^{\,8}\): take the product of their
specified full powers to obtain any eight coordinates.
No assertion about the whole abelianization of
\(D^+\) follows from this calculation.

\emph{The two section labels under this filling.}
The action on each two-point configuration fibre
fixes the two base points individually:
\((z_+,z_-)\mapsto(rz_+,rz_-)\).
Its fixed subgroup in \(\mathcal J_4\) is the kernel
of (F.19.13), and is therefore trivial.
Its continuous section classes have the same
two labels as in F.18.  Here are the inputs for
this smaller actual group.
The discrete extension
\(J_{4,\mathrm d}\rtimes C_2\) is the subgroup of
the six-point mapping class group generated by
the point-pushing kernel and the indicated involution.
Each order-two section exchanges the two old
pairs and fixes the two new labels.
Nielsen realization and the ordered quotient
homeomorphism identify it with the standard
involution by a pure six-point conjugator.
The double-coset comparison of those conjugators is
\(C_{4,\mathrm d}/C_{4,\mathrm d}^+=C_2\).
The label is the branch permutation of the
forgotten conjugator.
All their discrete fixed subgroups in
\(J_{4,\mathrm d}\) are trivial, by the discrete
version of (F.19.13).

The split free-by-free sequence (F.19.16) gives the
actual full topology and the two-row cohomology
comparison, just as in F.18.12.
Finite graphs for \(F_4\) and \(F_3\), with mapping
cylinders of representatives of their monodromy
homotopy equivalences, give a finite classifying
complex of dimension two.
The extension is residually finite as this actual
mapping class subgroup.  The finite extension
cohomology comparison and finite virtual
cohomological-type hypotheses are therefore the
same verified ones as in F.18.13--14.
The prime-order theorem \cite[Theorem~A]{R-BZ}
transports the two conjugacy classes and their
centralizers to the full completion.
Conjugacy by the extension can be changed to
conjugacy by its configuration kernel by removing
the quotient-generator power of the conjugator.
Thus the full \(H^1(C_2,\mathcal J_4)\) has exactly
these two branch labels.

The filling \(\mathcal J\to\mathcal J_4\) carries
the two full classes of F.18 to these two classes
with the same labels.
At the discrete level this follows by forgetting
the other two ordinary quotient orbits: it preserves
both named branch points and their permutation.
The finite branch characters are continuous by
the full-topology comparison, so the identical
label statement holds in the completions.
The induced map on these two cohomology sets is
therefore a bijection.

Let \(\delta\in D^+\) satisfy \(\delta R(\delta)=1\).
Its image in \(\mathcal J_4\) is the identity
cocycle.  The just-proved bijection forces its
\(\mathcal J\)-class to be the identity class.
By F.18 there is a unique \(b\in\mathcal J\) with
\(\delta=b^{-1}R(b)\).
Its image \(g(b)\) is fixed in \(\mathcal J_4\),
so is the identity.
Hence \(b\in D^+\).
Uniqueness also follows from
\(C_{D^+}(R)\subset C_{\mathcal J}(R)=1\).
The coboundary map is continuous; compactness of
\(D^+\) and the Hausdorff cocycle space make this
continuous bijection a homeomorphism.
This proves (F.19.4--5).

For \(z\in\mathcal Q^R\), choose \(k\in\pi^{-1}(z)\).
Then \(k^{-1}R(k)\) is a cocycle in \(D^+\).
Its unique corrector \(b\in D^+\) makes
\(kb^{-1}\) fixed, and does not change its image
under \(\pi\).
The kernel of \(\pi\) on the fixed subgroup is
\(C_{D^+}(R)=1\).
This proves (F.19.6), including its actual inverse
and continuity.

\emph{The specified source image inside the pullback.}
For \(u\in\mathcal E\), the element \(g\ell(u)\)
lies in \(C_6^\dagger\).
Its ordinary four-point image is
\(f_0\widehat\phi_*(u)\).
The exact source square (F.17.10--11) identifies
that with \(h\widehat\chi_*(q_0u)\).
The latter element also belongs to \(C_6^\dagger\)
by (F.19.14).
Injectivity in (F.19.13) proves the first equality
of (F.19.7).
For \(qw=u\), the same square gives
\(g\Phi^\dagger(w)=\widehat\chi_*(q_0u)\).
Thus \(g(b(w))=1\); its \(f\)-image is already
the identity by F.18.
This proves \(b(w)\in D^+\).
For \(qv=1\), both forgetting images of
\(\Phi^\dagger(v)\) are the identity, so
\(K^+\subset D^+\).
The right-coset comparison and its independence
are the actual equality
\(b(wv)=b(w)\Phi^\dagger(v)\).
The quotient map \(q^{-1}(\mathcal E)\to\mathcal E\)
is a compact surjection, so that equality and
continuity give the asserted continuous descended
map (F.19.8).
The inclusion \(D^+/K^+\to\mathcal J/K^+\)
is a closed embedding of compact coset spaces.
Hence its distinguished fibre is precisely
the earlier fibre (F.18.10).

The same source square puts \(\gamma(L)\) in
\(\mathcal Q\), and injectivity of
\(\widehat\phi_*\) makes \(\gamma\) an embedding.
Define \(\pi_L\) by the inverse of that embedding
after applying \(\pi\) to \(\mathcal T\).
It is onto by (F.19.2).
For \(z=\Phi^\dagger(w)\in H^\dagger\), the equality
\(\pi(z)=\gamma(qw)\) shows that \(q_H(z)=qw\)
is independent of \(w\).
This is a continuous surjection: it is induced
by a continuous map on the compact source.
Its kernel is exactly \(K^+\).
These facts prove every map and kernel in
(F.19.10).

Use the free based meridians in F.17.5.
The filling \(q\) kills exactly the sixteen new
generators and sends the eleven old finite
meridians to a free basis of \(L\).
Sending that latter basis back to those eleven
specified meridians defines a continuous map
\(i:L\to L^\dagger\) with \(qi=1\).
Thus \(j=\Phi^\dagger i\) splits both rows of
(F.19.10).
Conjugation by \(j(L)\) preserves \(D^+\), and
also preserves \(K^+\), the kernel in the upper
row.  The multiplication and inverse coordinates
in (F.19.11) are consequently homomorphisms
of the stated semidirect products and are
inverse continuous maps.
This constructs the common splitting directly
from the named free meridians.

Finally an element of \(\mathcal T\cap C^\dagger\)
has projection \(\gamma(u)\) for some \(u\in L\).
Its first component being fixed is exactly the
condition \(u\in\mathcal E\), by F.12.
The second component is already fixed by
(F.19.14).
Equation (F.19.6) gives a unique fixed lift of
that pair; (F.19.7) identifies it with \(\ell(u)\).
Conversely \(\ell(u)\) belongs to \(\mathcal T\)
for every \(u\in\mathcal E\).
This proves (F.19.12).

The new rational comparison checked by
\path{code/part_a/common_reflection_cross_kernel_f19.py}
is particularly explicit.
Writing \(w_j=v_{\zeta^j}\), \(p=a_+\), \(m=a_-\),
the normalized action sends
\[
 (w_0,w_1,w_2,w_3,w_4,p,m)\longmapsto
 \left(\frac{w_1}{w_2},w_1,\frac{w_1}{w_0},
       \frac{w_1}{w_4},\frac{w_1}{w_3},
       \frac{w_1}{p},\frac{w_1}{m}\right).
                                                               \tag{F.19.20}
\]
It fixes the retained curve
\((w_1,p,m)=(a^2,a,-a)\).
For \(\sigma=(0\ 2)(3\ 4)\) and \(z=p\) or \(m\),
its cross difference is exactly
\[
 R(w_j)-R(z)
       =-\frac{w_1}{w_{\sigma(j)}z}
                      \bigl(w_{\sigma(j)}-z\bigr).
                                                               \tag{F.19.21}
\]
The coefficient is nonzero in the relevant
configuration chart.  Thus the eight cross
divisors and their labels are preserved in
the stated four pairs.
The executable checks these identities, the
involution identity and the sixteen-column
winding matrix from the simple source collisions.
It passes 44 exact checks.
The fibre-product, completion and section arguments
above establish the group assertions; the calculation
does not identify \(K^+\) with \(D^+\).
\end{proof}
